% modul-3.tex - Modul 3: Datenverarbeitung & Dateien

\startcomponent modul-3
\environment ../pyt-layout
\environment ../pyt-env

\chapter{Datenverarbeitung \& Dateien}

{\it Dauer: 1 Tag (4 Lektionen à 50 Minuten) • Voraussetzung: Module 1-2 abgeschlossen}

\startlernziele
Nach diesem Modul können Sie:
\startitemize[n,packed]
    \item Dateien lesen und schreiben (TXT, CSV, JSON)
    \item Daten strukturiert verarbeiten
    \item Fehlerbehandlung mit try/except implementieren
    \item Mit externen Bibliotheken arbeiten
    \item API-Calls durchführen
    \item KI für Datenverarbeitung nutzen
\stopitemize
\stoplernziele

% =============================================================================
% VORBEREITUNG
% =============================================================================

\section{Vorbereitung}

{\it Zeitaufwand: 2-3 Stunden vor dem Präsenzunterricht}

\subsection{Leseauftrag: File I/O \& Datenformate}

\subsubsection{Dateiformate verstehen}

\starttabulate[|l|l|p(6cm)|]
\HL
\NC {\bf Format} \NC {\bf Verwendung} \NC {\bf Beispiel} \NC\NR
\HL
\NC TXT \NC Einfacher Text \NC Logs, Notizen \NC\NR
\NC CSV \NC Tabellendaten \NC Excel-Export, Datenbanken \NC\NR
\NC JSON \NC Strukturierte Daten \NC APIs, Konfiguration \NC\NR
\HL
\stoptabulate

\subsubsection{Experimente mit Dateien}

Erstellen Sie eine Testdatei und experimentieren Sie:

\startbashcode
echo "Zeile 1" > test.txt
echo "Zeile 2" >> test.txt
cat test.txt
\stopbashcode

\startpythoncode
# Python REPL
with open("test.txt", "r") as f:
    print(f.read())
\stoppythoncode

% =============================================================================
% PRAXIS
% =============================================================================

\section{Praxis}

{\it Präsenzunterricht: 4 Lektionen à 50 Minuten}

\subsection{Lektion 1: Dateien lesen \& schreiben}

{\it Dauer: 50 Minuten}

\subsubsection{Dateien öffnen}

\startpythoncode
# Ohne with (nicht empfohlen)
file = open("datei.txt", "r")
inhalt = file.read()
file.close()

# Mit with (empfohlen - schliesst automatisch)
with open("datei.txt", "r") as file:
    inhalt = file.read()
\stoppythoncode

\subsubsection{File Modes}

\starttabulate[|l|p(9cm)|]
\HL
\NC {\bf Mode} \NC {\bf Beschreibung} \NC\NR
\HL
\NC \code{"r"} \NC Read -- Lesen (Standard) \NC\NR
\NC \code{"w"} \NC Write -- Schreiben (überschreibt) \NC\NR
\NC \code{"a"} \NC Append -- Anhängen \NC\NR
\NC \code{"r+"} \NC Read + Write \NC\NR
\NC \code{"rb"} \NC Read Binary \NC\NR
\HL
\stoptabulate

\subsubsection{Lesen}

\startpythoncode
# Alles lesen
with open("datei.txt", "r") as f:
    inhalt = f.read()

# Zeile fuer Zeile
with open("datei.txt", "r") as f:
    for zeile in f:
        print(zeile.strip())

# Alle Zeilen als Liste
with open("datei.txt", "r") as f:
    zeilen = f.readlines()
\stoppythoncode

\subsubsection{Schreiben}

\startpythoncode
# Ueberschreiben
with open("datei.txt", "w") as f:
    f.write("Hallo Welt\n")
    f.write("Zweite Zeile\n")

# Anhaengen
with open("datei.txt", "a") as f:
    f.write("Neue Zeile\n")
\stoppythoncode

\subsubsection{Pfade mit pathlib}

\startpythoncode
from pathlib import Path

# Pfad erstellen
pfad = Path("daten/datei.txt")

# Pruefen ob existiert
if pfad.exists():
    print("Datei existiert")

# Verzeichnis erstellen
pfad.parent.mkdir(parents=True, exist_ok=True)
\stoppythoncode

\subsection{Lektion 2: CSV-Datenverarbeitung}

{\it Dauer: 50 Minuten}

\subsubsection{csv-Modul}

\startpythoncode
import csv

# Lesen
with open("daten.csv", "r") as f:
    reader = csv.reader(f)
    for row in reader:
        print(row)

# Mit DictReader (Spaltenname als Key)
with open("daten.csv", "r") as f:
    reader = csv.DictReader(f)
    for row in reader:
        print(row["name"], row["alter"])

# Schreiben
with open("output.csv", "w", newline="") as f:
    writer = csv.writer(f)
    writer.writerow(["Name", "Alter"])
    writer.writerow(["Anna", 25])
\stoppythoncode

\subsubsection{pandas Basics}

\startpythoncode
import pandas as pd

# CSV lesen
df = pd.read_csv("daten.csv")

# Anzeigen
print(df.head())      # Erste 5 Zeilen
print(df.info())      # Spalteninfos
print(df.describe())  # Statistiken

# Filtern
erwachsene = df[df["alter"] >= 18]

# Aggregieren
durchschnitt = df["alter"].mean()
summe = df["umsatz"].sum()

# Gruppieren
pro_kategorie = df.groupby("kategorie")["umsatz"].sum()

# Schreiben
df.to_csv("output.csv", index=False)
\stoppythoncode


\subsection{Lektion 3: JSON \& APIs}

{\it Dauer: 50 Minuten}

\subsubsection{JSON verarbeiten}

\startpythoncode
import json

# JSON lesen
with open("config.json", "r") as f:
    daten = json.load(f)

# Auf Daten zugreifen
print(daten["name"])
print(daten["einstellungen"]["theme"])

# JSON schreiben
daten = {
    "name": "Meine App",
    "version": "1.0",
    "debug": True
}

with open("config.json", "w") as f:
    json.dump(daten, f, indent=2)

# String zu JSON
json_string = '{"name": "Anna", "alter": 25}'
person = json.loads(json_string)

# JSON zu String
person_string = json.dumps(person, indent=2)
\stoppythoncode

\subsubsection{API-Calls mit requests}

\startpythoncode
import requests

# GET-Request
response = requests.get("https://api.example.com/data")

if response.status_code == 200:
    daten = response.json()
    print(daten)
else:
    print(f"Fehler: {response.status_code}")

# Mit Parametern
params = {"city": "Zurich", "units": "metric"}
response = requests.get("https://api.weather.com", params=params)

# POST-Request
daten = {"name": "Anna", "email": "anna@example.com"}
response = requests.post("https://api.example.com/users", json=daten)
\stoppythoncode

\startwarnbox
{\bf API-Keys sicher aufbewahren!}

Niemals API-Keys in den Code schreiben. Verwenden Sie Umgebungsvariablen:

\startpythoncode
import os
api_key = os.getenv("API_KEY")
\stoppythoncode
\stopwarnbox

\subsection{Lektion 4: Fehlerbehandlung \& Logging}

{\it Dauer: 50 Minuten}

\subsubsection{try/except}

\startpythoncode
try:
    with open("datei.txt", "r") as f:
        inhalt = f.read()
except FileNotFoundError:
    print("Datei nicht gefunden!")
except PermissionError:
    print("Keine Berechtigung!")
except Exception as e:
    print(f"Unbekannter Fehler: {e}")
finally:
    print("Wird immer ausgefuehrt")
\stoppythoncode

\subsubsection{Eigene Exceptions}

\startpythoncode
class ValidationError(Exception):
    """Fehler bei Datenvalidierung."""
    pass

def validiere_alter(alter: int) -> None:
    if alter < 0:
        raise ValidationError("Alter kann nicht negativ sein")
    if alter > 150:
        raise ValidationError("Alter unrealistisch hoch")

try:
    validiere_alter(-5)
except ValidationError as e:
    print(f"Validierungsfehler: {e}")
\stoppythoncode

\subsubsection{Logging}

\startpythoncode
import logging

# Konfiguration
logging.basicConfig(
    level=logging.INFO,
    format="%(asctime)s - %(levelname)s - %(message)s",
    filename="app.log"
)

# Verwenden
logging.debug("Debug-Information")
logging.info("Programm gestartet")
logging.warning("Warnung!")
logging.error("Fehler aufgetreten")
logging.critical("Kritischer Fehler!")
\stoppythoncode

% =============================================================================
% ÜBUNGEN
% =============================================================================

\section{Übungen}

\startaufgabenbox
{\bf Übung 1: Log-Datei-Analyse} {\it (15 Min.)}

Analysieren Sie eine Log-Datei:
\startitemize[packed]
    \item Anzahl ERROR, WARNING, INFO zählen
    \item Häufigste Fehlermeldung finden
    \item Zusammenfassung ausgeben
\stopitemize
\stopaufgabenbox

\blank[medium]

\startaufgabenbox
{\bf Übung 2: CSV-Datenbereinigung} {\it (15 Min.)}

CSV-Datei bereinigen:
\startitemize[packed]
    \item Leere Zeilen entfernen
    \item Duplikate entfernen
    \item Fehlende Werte behandeln
    \item Bereinigte Datei speichern
\stopitemize
\stopaufgabenbox

\blank[medium]

\startaufgabenbox
{\bf Übung 3: JSON-Konfiguration} {\it (15 Min.)}

Konfigurationssystem erstellen:
\startitemize[packed]
    \item config.json lesen
    \item Default-Werte bei fehlenden Keys
    \item Änderungen speichern
    \item Fehlerbehandlung
\stopitemize
\stopaufgabenbox

\blank[medium]

\startaufgabenbox
{\bf Übung 4: API-Integration} {\it (20 Min.)}

Öffentliche API nutzen (z.B. OpenWeatherMap):
\startitemize[packed]
    \item API-Call durchführen
    \item Daten parsen
    \item Formatiert ausgeben
    \item Fehler behandeln
\stopitemize
\stopaufgabenbox

% =============================================================================
% NACHBEARBEITUNG
% =============================================================================

\section{Nachbearbeitung}

{\it Hausaufgaben: 4-6 Stunden vor Modul 4}

\subsection{Aufgabe 1: Daten-Pipeline}

Erstellen Sie eine Pipeline:
\startitemize[n,packed]
    \item CSV-Datei einlesen
    \item Daten validieren und bereinigen
    \item Als JSON exportieren
    \item Logging für jeden Schritt
\stopitemize

\subsection{Aufgabe 2: CSV-zu-JSON-Konverter}

CLI-Tool erstellen:

\startbashcode
python konverter.py eingabe.csv ausgabe.json
\stopbashcode

\startitemize[packed]
    \item Beliebige CSV-Dateien konvertieren
    \item Datentypen automatisch erkennen
    \item Fortschritt anzeigen
\stopitemize

\subsection{Aufgabe 3: Web-Scraper}

Einfachen Scraper implementieren:
\startitemize[packed]
    \item Webseite abrufen mit requests
    \item Daten extrahieren
    \item Als CSV speichern
    \item robots.txt beachten!
\stopitemize

\subsection{Aufgabe 4: Datenanalyse-Projekt}

Öffentlichen Datensatz analysieren:
\startitemize[packed]
    \item Daten herunterladen (z.B. von Kaggle)
    \item Mit pandas analysieren
    \item Mindestens 5 Erkenntnisse dokumentieren
    \item Ergebnisse visualisieren (optional: matplotlib)
\stopitemize

% =============================================================================
% MATERIALIEN
% =============================================================================

\section{Materialien}

\subsection{File I/O Cheat Sheet}

\starttabulate[|l|p(8cm)|]
\HL
\NC {\bf Aktion} \NC {\bf Code} \NC\NR
\HL
\NC Lesen \NC \code{with open("f.txt", "r") as f: f.read()} \NC\NR
\NC Schreiben \NC \code{with open("f.txt", "w") as f: f.write(text)} \NC\NR
\NC Anhängen \NC \code{with open("f.txt", "a") as f: f.write(text)} \NC\NR
\NC CSV lesen \NC \code{pd.read_csv("f.csv")} \NC\NR
\NC JSON lesen \NC \code{json.load(f)} \NC\NR
\NC API-Call \NC \code{requests.get(url).json()} \NC\NR
\HL
\stoptabulate

\subsection{Error Handling Guide}

\startpythoncode
# Muster fuer robuste Dateiverarbeitung
def verarbeite_datei(pfad: str) -> dict:
    """Verarbeitet Datei mit Fehlerbehandlung."""
    try:
        with open(pfad, "r") as f:
            daten = json.load(f)
        logging.info(f"Datei {pfad} geladen")
        return daten
    except FileNotFoundError:
        logging.error(f"Datei nicht gefunden: {pfad}")
        return {}
    except json.JSONDecodeError as e:
        logging.error(f"JSON-Fehler: {e}")
        return {}
\stoppythoncode

\blank[big]

\starttippbox
{\bf Erfolgskriterien für Modul 3}

Sie haben das Modul erfolgreich abgeschlossen, wenn Sie:
\startitemize[packed]
    \item Dateien sicher lesen und schreiben können
    \item CSV-Daten mit csv oder pandas verarbeiten können
    \item JSON parsen und erstellen können
    \item API-Calls durchführen können
    \item Robuste Fehlerbehandlung implementieren können
\stopitemize
\stoptippbox

\stopcomponent

